\newcommand{\VtwoN}{9254}
\newcommand{\VtwoAdults}{5569}
\newcommand{\VtwoOlder}{2150}
\newcommand{\VtwoDf}{15}
\newcommand{\VtwoWeighted}{29{,}1516}
\newcommand{\VtwoUnweighted}{38{,}6066}
\newcommand{\VtwoSE}{1{,}5529}
\newcommand{\VtwoLower}{25{,}8416}
\newcommand{\VtwoUpper}{32{,}4615}
\subsection{Sayısal uygulama: dosya oranı ve hedef nüfus oranı}
\label{sec:v2-nhanes-sonuclar}

Kaynakta \VtwoN{} kişi kaydı vardır; bunların \VtwoAdults{} tanesi
20 yaş ve üzerinde, \VtwoOlder{} tanesi 60 yaş ve üzerindedir.
Seçili alanlarda eksiksizlik, pozitif görüşme ağırlıkları ve 15
tabaka--30 PSU yapısı kontrol edilmiştir. Diğer sütunların eksikleri
nedeniyle kayıt silinmemiştir. Dosyanın sayısal kabulü, kamu dosyasının
gizli hane seçim çerçevesini içerdiği anlamına gelmez.

\begin{center}
\begin{tabular}{@{}lr@{}}
\toprule
Özet & Sonuç \\
\midrule
Ağırlıksız yetişkin kayıt oranı & \%\VtwoUnweighted \\
Ağırlıklı alt evren oranı & \%\VtwoWeighted \\
Taylor standart hatası & \VtwoSE{} yüzde puan \\
Tasarım serbestlik derecesi & \VtwoDf \\
Yaklaşık \%95 güven aralığı & [\%\VtwoLower; \%\VtwoUpper] \\
\bottomrule
\end{tabular}
\end{center}

Ağırlıklı pay yaklaşık 69,596 milyon, payda 238,738 milyondur;
bunlar dosyadaki kişi sayısı değil, görüşme ağırlıklarıyla elde edilen
nüfus toplamı tahminleridir. Ağırlıksız oran yanıtlayıcıların bileşimini,
ağırlıklı oran hedef alt nüfusu özetler. Farklı olmaları hesap çelişkisi
veya iki bağımsız araştırma kanıtı değildir.

\textbf{İzlenebilir varyans hesabı.} $B=\sum_iw_id_i$ ve
$z_i=w_id_i(y_i-\widehat p_w)/B$ olsun. Tabaka $h$ içindeki PSU $j$
için $Z_{hj}$ kişi katkılarının toplamıdır. Nihai PSU yaklaşımıyla,
sonlu evren düzeltmesi kullanmadan,
\[
\widehat{\Var}(\widehat p_w)=
\sum_h\frac{m_h}{m_h-1}\sum_j(Z_{hj}-\overline Z_h)^2
\]
hesaplanır. Bütün kayıtlar tasarımda kalır; çocukların alt evren
katkıları sıfırdır. Bu veride yetişkinler bütün 15 tabaka ve 30 PSU'da
bulunduğundan serbestlik derecesi $30-15=15$ alınır
\citep{nchsvariance}. PSU kodları tabaka koduyla birlikte kullanılır.

Aralık $\widehat p_w\pm t_{0{,}975;15}\widehat{\SE}$ biçiminde,
yaklaşık t-Wald öğretim aralığıdır. NCHS'nin bütün oran yayımlama
ölçütlerinin uygulandığı iddia edilmez. Kapsam veya yanıtsızlık
yanlılığı yalnız standart hata hesabıyla ölçülmüş olmaz.

\begin{researchbox}
\textbf{Örnek rapor.} NHANES 2017--2018 görüşme ağırlıkları ve
maskelenmiş varyans tasarımıyla, hedefteki kurum dışı sivil yetişkin
nüfusta 60 yaş ve üzeri oranı yaklaşık \%29,15 tahmin edilmiştir
(Taylor standart hatası 1,55 yüzde puan; yaklaşık \%95 t aralığı
\%25,84--\%32,46). Bu tahmin bugünkü nüfusa, Türkiye'ye veya kurum
nüfusuna doğrudan taşınamaz; nedensel etkiyi ölçmez. CDC'nin belirli
bir resmî tahmininin tekrarı olarak sunulmaz.
\end{researchbox}

\begin{softwarebox}
Tam duyarlıklı sonuçlar \path{companion/vakalar/v2-nhanes/results.json},
30 PSU'nun katkıları aynı klasörde \texttt{psu.csv} içindedir.
\texttt{export.py} çıktıları üretir; \texttt{verify\_delivery.py}
dosya bütünlüğünü, sayısal eşleşmeyi ve kitap bağlantısını denetler.
İkinci hesap pay/payda toplamlarının kovaryansından delta yöntemiyle
aynı varyansa ulaşmıştır. XPT okuyucusu ve t kritik değeri ortaktır;
R/SPSS veya bağımsız bir anket yazılımı çalıştırılmış sayılmaz.
Sayısal eşleşme, örnekleme kapsamının veya yanıt sürecinin kusursuzluğunu
kanıtlamaz.
\end{softwarebox}